% Compile with XeLaTeX. Content converted from semestrova-kontrolna-stem-5-klas.docx.
\documentclass[11pt,a4paper]{article}
\usepackage{fontspec}
\usepackage[ukrainian]{babel}
\usepackage[margin=18mm]{geometry}
\usepackage{array,tabularx}
\setmainfont{Arial}
\setlength{\parindent}{0pt}
\setlength{\parskip}{6pt}
\renewcommand{\arraystretch}{1.65}
\begin{document}
\begin{center}\Large\textbf{Відповіді для вчителя}\par STEM, 5 клас\end{center}

\section*{Варіант 1}

Тести: 1 - Б; 2 - А; 3 - В.\par

Теоретичне питання, орієнтовна відповідь: Зображення, символи та знаки допомагають людині передавати інформацію. За допомогою зображень можна показати предмет, явище або подію. Символи передають ідею або значення у короткій формі. Знаки допомагають орієнтуватися в просторі, навчанні та спілкуванні. Наприклад, дорожні знаки попереджають, а герб або емблема показують належність. Тому зображення і знаки важливі в повсякденному житті.\par

\begin{tabularx}{\linewidth}{|>{\raggedright\arraybackslash}X|>{\raggedright\arraybackslash}X|>{\raggedright\arraybackslash}X|}\hline

\textbf{Космос} & \textbf{Графіка} & \textbf{Символи} \\ \hline

телескоп & фотографія & герб \\ \hline

орбіта & анімація & емблема \\ \hline

ракета & малюнок & піктограма \\ \hline

планета & колаж & печатка \\ \hline

\end{tabularx}\par

\begin{tabularx}{\linewidth}{|>{\raggedright\arraybackslash}X|>{\raggedright\arraybackslash}X|>{\raggedright\arraybackslash}X|}\hline

\textbf{Об'єкт} & \textbf{Для чого використовується} & \textbf{До якого модуля належить} \\ \hline

Телескоп & Для спостереження за небесними тілами & Космос \\ \hline

Герб & Для позначення належності, історії, традицій & Символи \\ \hline

Анімація & Для створення рухомих зображень & Графіка \\ \hline

Ракета & Для польотів у космос & Космос \\ \hline

Піктограма & Для швидкого передавання короткої інформації & Символи \\ \hline

Фотографія & Для фіксації та передавання зображень & Графіка \\ \hline

Печатка & Для позначення або підтвердження & Символи \\ \hline

Орбіта & Для позначення шляху руху тіла в космосі & Космос \\ \hline

\end{tabularx}\par

\newpage

\section*{Варіант 2}

Тести: 1 - А; 2 - А; 3 - В.\par

Теоретичне питання, орієнтовна відповідь: Знакова система - це сукупність знаків, за допомогою яких люди передають інформацію. Знаки бувають природні та штучні. Природні знаки існують у природі, наприклад темні хмари можуть означати дощ. Штучні знаки створює людина: букви, цифри, дорожні знаки, емблеми. Вони допомагають швидко повідомляти важливу інформацію. Без знакових систем було б складніше навчатися, спілкуватися й орієнтуватися.\par

\begin{tabularx}{\linewidth}{|>{\raggedright\arraybackslash}X|>{\raggedright\arraybackslash}X|>{\raggedright\arraybackslash}X|}\hline

\textbf{Небезпечний фактор} & \textbf{Чим загрожує} & \textbf{Спосіб захисту} \\ \hline

Холод & Переохолодженням & Теплозахисний скафандр \\ \hline

Космічне сміття & Ударами та пошкодженнями & Міцний захисний шар, екранування \\ \hline

Радіація & Шкодою для здоров'я & Спеціальні матеріали захисту, укриття \\ \hline

\end{tabularx}\par

Практичне завдання 2, зразок плану: 1. Обрати 3-4 космічні професії. 2. Зібрати коротку інформацію про них. 3. Добрати ілюстрації або піктограми. 4. Розмістити інформацію на плакаті. 5. Виділити потрібні навички для кожної професії. 6. Оформити заголовок і висновок.\par

\newpage

\section*{Варіант-зразок}

Тести: 1 - Б; 2 - А; 3 - В.\par

Теоретичне питання, орієнтовна відповідь: Люди використовують знаки й зображення, щоб швидко передавати інформацію. Малюнок або фотографія допомагає побачити предмет чи явище. Знаки й символи можуть попереджати, пояснювати або позначати щось важливе. Наприклад, дорожній знак допомагає орієнтуватися на дорозі. Герб або емблема показують належність до міста, школи чи команди. Тому знаки та зображення є важливою частиною нашого життя.\par

\begin{tabularx}{\linewidth}{|>{\raggedright\arraybackslash}X|>{\raggedright\arraybackslash}X|>{\raggedright\arraybackslash}X|}\hline

\textbf{Космос} & \textbf{Графіка} & \textbf{Символи} \\ \hline

ракета & малюнок & печатка \\ \hline

планета & фотографія & знак \\ \hline

телескоп & анімація & емблема \\ \hline

орбіта &  & герб \\ \hline

 &  & піктограма \\ \hline

\end{tabularx}\par

\begin{tabularx}{\linewidth}{|>{\raggedright\arraybackslash}X|>{\raggedright\arraybackslash}X|>{\raggedright\arraybackslash}X|}\hline

\textbf{Приклад} & \textbf{До якої групи належить} & \textbf{Коротке пояснення} \\ \hline

Телескоп & Космос & Прилад для спостереження за небесними тілами \\ \hline

Герб & Символи & Передає належність і має символічне значення \\ \hline

Анімація & Графіка & Створює рухоме зображення \\ \hline

Планета & Космос & Небесне тіло Сонячної системи або іншої системи \\ \hline

Піктограма & Символи & Умовний знак для швидкого повідомлення \\ \hline

Фотографія & Графіка & Фіксує й передає зображення \\ \hline

\end{tabularx}\par
\end{document}
